\documentclass[a4paper]{arrowhead}

\usepackage[yyyymmdd]{datetime}
\usepackage{etoolbox}
\usepackage[T1]{fontenc}
\usepackage[utf8]{inputenc}
\usepackage{multirow}
\usepackage{longtable}
\usepackage{amssymb}
\usepackage{float}
\usepackage{tikz}
\usetikzlibrary{positioning,arrows.meta,fit,calc,shapes.misc,shapes.geometric}

\renewcommand{\dateseparator}{-}
\setlength{\emergencystretch}{3em}
\lstset{basicstyle=\ttfamily\footnotesize,breaklines=true,columns=fullflexible,frame=single,rulecolor=\color{gray!50},backgroundcolor=\color{gray!6}}
\tikzset{
  ahsys/.style={draw=ArrowheadBlue,very thick,fill=ArrowheadBlue!6,rounded corners=2pt,minimum width=3.2cm,minimum height=1.1cm,align=center,font=\small\bfseries},
  ahext/.style={draw=ArrowheadGrey,thick,dashed,fill=gray!5,rounded corners=2pt,minimum width=3.2cm,minimum height=1.0cm,align=center,font=\small},
  ahsrv/.style={draw=ArrowheadPurple,thick,fill=ArrowheadPurple!8,rounded corners=6pt,minimum height=0.55cm,inner xsep=4pt,align=center,font=\scriptsize\ttfamily},
  ahcsm/.style={draw=ArrowheadTeal,thick,fill=ArrowheadTeal!8,rounded corners=6pt,minimum height=0.55cm,inner xsep=4pt,align=center,font=\scriptsize\ttfamily},
  ahdb/.style={draw=ArrowheadOrange,thick,fill=ArrowheadOrange!8,cylinder,shape border rotate=90,aspect=0.25,minimum height=0.9cm,align=center,font=\scriptsize},
  ahflow/.style={-{Stealth[length=2.2mm]},thick,ArrowheadGrey},
  ahlabel/.style={font=\scriptsize,fill=white,inner sep=1pt,align=center},
}

%% Special references
\newcommand{\mref}[1]{{\textcolor{ArrowheadPurple}{\hyperref[sec:model:#1]{#1}}}}
\newcommand{\oref}[1]{{\textcolor{ArrowheadBlue}{\hyperref[sec:operations:#1]{#1}}}}
\newcommand{\pdef}[1]{{\textcolor{ArrowheadGrey}{#1\label{sec:model:primitives:#1}\label{sec:model:primitives:#1s}\label{sec:model:primitives:#1es}}}}
\newcommand{\pref}[1]{{\textcolor{ArrowheadGrey}{\hyperref[sec:model:primitives:#1]{#1}}}}

\newrobustcmd\osubsection[3]{
  \addtocounter{subsection}{1}
  \addcontentsline{toc}{subsection}{\protect\numberline{\thesubsection}operation \textcolor{ArrowheadBlue}{#1}}
  \subsection*{
    \thesubsection\quad
    operation
    \textcolor{ArrowheadBlue}{#1}
    (\notblank{#2}{\mref{#2}}{})
    \notblank{#3}{: \mref{#3}}{}
  }
  \label{sec:operations:#1}
}
\newrobustcmd\msubsection[2]{
  \addtocounter{subsection}{1}
  \addcontentsline{toc}{subsection}{\protect\numberline{\thesubsection}#1 \textcolor{ArrowheadPurple}{#2}}
  \subsection*{\thesubsection\quad#1 \textcolor{ArrowheadPurple}{#2}}
  \label{sec:model:#2} \label{sec:model:#2s} \label{sec:model:#2es}
}

\begin{document}

%% Arrowhead Document Properties
\ArrowheadTitle{workflowConfiguration}
\ArrowheadServiceID{workflowConfiguration}
\ArrowheadType{Service Description}
\ArrowheadTypeShort{SD}
\ArrowheadVersion{5.4.0}
\ArrowheadDate{2026-11-24}
\ArrowheadAuthor{Documentation maintenance}
\ArrowheadStatus{REVIEW DRAFT}
\ArrowheadContact{Repository maintainers}
\ArrowheadFooter{\href{www.arrowhead.eu}{www.arrowhead.eu}}
\ArrowheadSetup
%%

%% Front Page
\begin{center}
  \vspace*{1cm}
  \huge{\arrowtitle}

  \vspace*{0.2cm}
  \LARGE{\arrowtype}
  \vspace*{1cm}

  \Large{Service ID: \textit{``\arrowid''}}
  \vspace*{\fill}

  

  \vspace*{1cm}
  \vspace*{\fill}

  % Front Page Abstract
  \begin{abstract}
    This document describes the \texttt{workflow\allowbreak{}Configuration} service of the Eclipse Arrowhead workflow support systems, which lists configured local clouds, their microservices and executors, and validates and stores edited workflows. It is provided by the WorkflowConfigurator system, service version 1.0.0, documentation edition 5.4.0.
  \end{abstract}

  \vspace*{1cm}
\end{center}

\newpage
%%

%% Table of Contents
\tableofcontents
\newpage
%%

\section{Overview}
\label{sec:overview}
The \texttt{workflow\allowbreak{}Configuration} service is what the graphical workflow editor of the
WorkflowConfigurator uses, offered as an Arrowhead service so that other engineering tools can use it as well. It
lists the local clouds a workflow can be designed for (the running cloud, LocalCloudConfigurator designs and the member
clouds of System of LocalClouds designs), the microservices offered in a cloud with their operations, and the
WorkflowExecutors with their capabilities. It forwards validation and storage of edited workflows to the
WorkflowManager, where recipes are stored and nowhere else.

The rest of this document is organized as follows.
In Section~\ref{sec:operations}, we describe the abstract message operations provided by the service.
In Section~\ref{sec:model}, we end the document by presenting the data types used by the mentioned operations.

\subsection{Significant Prior Art}
Garcia Represa's workflow selection and implementation phases~\cite{represa2022}; the
LocalCloudConfigurator and SoLocalCloudConfigurator configuration services of this workspace.

\subsection{How This Service Is Meant to Be Used}
The editor first calls \oref{list-clouds} and lets the operator pick a cloud, then \oref{list-services} and
\oref{list-executors} to populate its palette. While drawing it calls \oref{validate-recipe}; \oref{store-recipe}
registers the workflow in the WorkflowManager when it validates. \oref{list-recipes} and \oref{get-recipe} open stored
workflows.

\subsection{Important Delimitations and Dependencies}
\label{sec:delimitations}
\begin{itemize}
\item Every source is optional; an unavailable source is reported in \texttt{sources} and the others are still listed.
\item Services of another cloud are listed only when that cloud is the live one; its ServiceRegistry is not reachable from here.
\item Without a \texttt{service\allowbreak{}Registry\allowbreak{}Management} grant the configurator falls back to discovery, which under the RESTRICTED discovery policy hides services offered without \texttt{unrestricted\allowbreak{}Discovery}.
\item The configurator stores nothing itself; validation and storage are the WorkflowManager's.
\end{itemize}

\paragraph{Authorization.} Each operation is authorized separately at the Arrowhead connector. The browser console reaches the same operations on a separate loopback connector with its own password login (Section on the console).

\newpage
\section{Service Interface}
\label{sec:operations}
This section describes the interface of the \texttt{workflow\allowbreak{}Configuration} service. Each subsection names an abstract operation,
an input type and an output type, in that order. The input type is named inside parentheses, while the output type is
preceded by a colon. Input and output types are only denoted when accepted or returned, respectively, by the operation
in question. All abstract data types named in this section are defined in Section~\ref{sec:model}. The concrete
interface (HTTP/TLS/JSON) is specified in the IDD of the same name.

\begin{figure}[ht!]
  \centering
  \resizebox{0.92\textwidth}{!}{
\begin{tikzpicture}
\node[ahsys] (prov) at (4.6,0) {WorkflowConfigurator\\[-1pt]{\normalfont\scriptsize $\ll$provider$\gg$}};
\node[ahsrv] (srv) at (1.6,0) {workflowConfiguration};
\draw[thick,ArrowheadPurple] (srv.east) -- (prov.west);
\node[ahext] (c0) at (-3.0,0.60) {Browser console (operator)\\[-1pt]{\scriptsize $\ll$consumer$\gg$}};
\draw[ahflow] (c0.east) -- (srv.west);
\node[ahext] (c1) at (-3.0,-0.60) {Engineering tool\\[-1pt]{\scriptsize $\ll$consumer$\gg$}};
\draw[ahflow] (c1.east) -- (srv.west);
\end{tikzpicture}
  }
  \caption{SysML block view of the \texttt{workflowConfiguration} service: provider and typical consumers.}
  \label{fig:SD}
\end{figure}

The following interface operations are available.

{\small
\begin{longtable}{|p{4.4cm}|p{5.6cm}|p{5.8cm}|}

\hline
\rowcolor{gray!33} Operation & Input & Output \\ \hline
\endhead
\oref{list-clouds} & -- & \mref{CloudList} \\ \hline
\oref{list-services} & \mref{CloudRef} & \mref{ServiceList} \\ \hline
\oref{list-executors} & -- & \mref{ExecutorCapabilityList} \\ \hline
\oref{list-recipes} & \mref{RecipeFilter} & \mref{RecipeList} \\ \hline
\oref{get-recipe} & \mref{RecipeKey} & \mref{RecipeRecord} \\ \hline
\oref{validate-recipe} & \mref{Recipe} & \mref{ValidationReport} \\ \hline
\oref{store-recipe} & \mref{Recipe} & \mref{StoreResult} \\ \hline
\end{longtable}
}

\osubsection{list-clouds}{}{CloudList}
Lists the live cloud, the LocalCloudConfigurator designs and the SoLC member clouds.
\osubsection{list-services}{CloudRef}{ServiceList}
Lists the microservices of a cloud from the ServiceRegistry, without \texttt{sysMonitor} and \texttt{general\allowbreak{}Management}.
\osubsection{list-executors}{}{ExecutorCapabilityList}
Lists the WorkflowExecutors of the live cloud with their capabilities.
\osubsection{list-recipes}{RecipeFilter}{RecipeList}
Forwards \texttt{workflow\allowbreak{}Recipe.\allowbreak{}query} with \texttt{summary=true}: the entries carry no \texttt{recipe} document (a full list outgrew the HTTP client's 256~KiB response buffer); \oref{get-recipe} reads the chosen version.
\osubsection{get-recipe}{RecipeKey}{RecipeRecord}
Forwards \texttt{workflow\allowbreak{}Recipe.\allowbreak{}get}.
\osubsection{validate-recipe}{Recipe}{ValidationReport}
Forwards \texttt{workflow\allowbreak{}Recipe.\allowbreak{}validate}.
\osubsection{store-recipe}{Recipe}{StoreResult}
Validates with the WorkflowManager and registers the workflow when the verdict is ACCEPT.

\newpage
\section{Information Model}
\label{sec:model}
Here, all data objects that can be part of the \texttt{workflow\allowbreak{}Configuration} service are listed in alphabetic order. Each
subsection, which describes one type of object, begins with the \textit{struct} keyword (a collection of named fields,
each with its own data type), the \textit{enum} keyword (a closed set of values) or the \textit{alias} keyword. As a
complement to the explicitly defined types in this section, there is also a list of implicit primitive types in
Section~\ref{sec:model:primitives}.

\msubsection{struct}{Action}
An action $\alpha$: a group of steps run in one work centre.

{\small
\begin{longtable}{|p{3.9cm}|p{4.4cm}|p{7.5cm}|}

\hline
\rowcolor{gray!33} Field & Type & Description \\ \hline
\endhead
\texttt{name} & \pref{Name} & Action name. \\ \hline
\texttt{next} & \pref{Name} & The action that becomes active when this one finishes; absent for the last action. Optional. \\ \hline
\texttt{workCentre} & \pref{String} & Informative work-centre identifier. Optional. \\ \hline
\texttt{dispatch} & \mref{DispatchMode} & Dispatch mode of this action's steps. Optional. \\ \hline
\texttt{steps} & \pref{List}$<$\mref{Step}$>$ & The steps $S_\alpha$. \\ \hline
\end{longtable}
}

\msubsection{struct}{CloudEntry}
One configured local cloud.

{\small
\begin{longtable}{|p{3.9cm}|p{4.4cm}|p{7.5cm}|}

\hline
\rowcolor{gray!33} Field & Type & Description \\ \hline
\endhead
\texttt{id} & \pref{String} & \texttt{live}, \texttt{lcc:<design>} or \texttt{solc:<design>:<cloud>}. \\ \hline
\texttt{name} & \pref{String} & Display name. \\ \hline
\texttt{cloud} & \pref{String} & Local cloud name. \\ \hline
\texttt{source} & \pref{SystemName} & ServiceRegistry, LocalCloudConfigurator or SoLocalCloudConfigurator. \\ \hline
\texttt{live} & \pref{Boolean} & Whether this cloud's ServiceRegistry is the one reachable from here. \\ \hline
\texttt{description} & \pref{String} & Description. \\ \hline
\texttt{systems} & \pref{List}$<$\mref{CloudSystem}$>$ & Systems of the cloud. \\ \hline
\texttt{revision} & \pref{Integer} & Design revision (LocalCloudConfigurator designs). Optional. \\ \hline
\texttt{updatedAt} & \pref{DateTime} & Design update time. Optional. \\ \hline
\end{longtable}
}

\msubsection{struct}{CloudList}
The local clouds the editor can design workflows for.

{\small
\begin{longtable}{|p{3.9cm}|p{4.4cm}|p{7.5cm}|}

\hline
\rowcolor{gray!33} Field & Type & Description \\ \hline
\endhead
\texttt{entries} & \pref{List}$<$\mref{CloudEntry}$>$ & Clouds. \\ \hline
\texttt{sources} & \pref{List}$<$\mref{SourceState}$>$ & Availability of each source; an unavailable source does not hide the others. \\ \hline
\texttt{note} & \pref{String} & Explanation when a fallback was used. Optional. \\ \hline
\end{longtable}
}

\msubsection{struct}{CloudRef}
Identifies a cloud of \oref{list-clouds}.

{\small
\begin{longtable}{|p{3.9cm}|p{4.4cm}|p{7.5cm}|}

\hline
\rowcolor{gray!33} Field & Type & Description \\ \hline
\endhead
\texttt{id} & \pref{String} & Cloud identifier. \\ \hline
\end{longtable}
}

\msubsection{struct}{CloudSystem}
A system of a cloud.

{\small
\begin{longtable}{|p{3.9cm}|p{4.4cm}|p{7.5cm}|}

\hline
\rowcolor{gray!33} Field & Type & Description \\ \hline
\endhead
\texttt{name} & \pref{SystemName} & System. \\ \hline
\texttt{running} & \pref{Boolean} & Whether it is registered in the live cloud. \\ \hline
\texttt{kind} & \pref{String} & Component kind in a design. Optional. \\ \hline
\end{longtable}
}

\msubsection{enum}{DispatchMode}
How a step reaches an executor.

{\small
\begin{longtable}{|p{3.9cm}|p{11.9cm}|}

\hline
\rowcolor{gray!33} Value & Meaning \\ \hline
\endhead
\texttt{push} & The WorkflowManager selects an executor and calls \texttt{workflow\allowbreak{}Execution.\allowbreak{}start} (default). \\ \hline
\texttt{pull} & The WorkflowManager publishes the work item; an executor claims it with \texttt{workflow\allowbreak{}Work\allowbreak{}Item.\allowbreak{}claim}. \\ \hline
\texttt{event} & Reserved for EventHandler notification; currently handled as \texttt{pull}. \\ \hline
\end{longtable}
}

\msubsection{struct}{ExecutorCapability}
An executor's declared capabilities.

{\small
\begin{longtable}{|p{3.9cm}|p{4.4cm}|p{7.5cm}|}

\hline
\rowcolor{gray!33} Field & Type & Description \\ \hline
\endhead
\texttt{name} & \pref{SystemName} & Executor. \\ \hline
\texttt{capabilities} & \pref{List}$<$\pref{String}$>$ & Capabilities. \\ \hline
\texttt{services} & \pref{List}$<$\pref{ServiceDefinition}$>$ & Runnable services. \\ \hline
\texttt{mode} & \mref{DispatchMode} & Mode. \\ \hline
\texttt{workstation} & \pref{String} & Workstation. Optional. \\ \hline
\end{longtable}
}

\msubsection{struct}{ExecutorCapabilityList}
Executors registered in the live cloud.

{\small
\begin{longtable}{|p{3.9cm}|p{4.4cm}|p{7.5cm}|}

\hline
\rowcolor{gray!33} Field & Type & Description \\ \hline
\endhead
\texttt{entries} & \pref{List}$<$\mref{ExecutorCapability}$>$ & Executors. \\ \hline
\end{longtable}
}

\msubsection{struct}{ExecutorRequirement}
$\eta(s)$: constraints on the executor that runs a step.

{\small
\begin{longtable}{|p{3.9cm}|p{4.4cm}|p{7.5cm}|}

\hline
\rowcolor{gray!33} Field & Type & Description \\ \hline
\endhead
\texttt{capability} & \pref{String} & A capability the executor must declare in \texttt{workflow\allowbreak{}Capabilities}. Optional. \\ \hline
\texttt{workstation} & \pref{String} & The executor's workstation. Optional. \\ \hline
\texttt{system} & \pref{SystemName} & A named executor system. Optional. \\ \hline
\end{longtable}
}

\msubsection{struct}{Finding}
One validator finding.

{\small
\begin{longtable}{|p{3.9cm}|p{4.4cm}|p{7.5cm}|}

\hline
\rowcolor{gray!33} Field & Type & Description \\ \hline
\endhead
\texttt{code} & \pref{String} & Check code, e.g. \texttt{D03}; the codes are listed in the \texttt{workflow\allowbreak{}Recipe} SD. \\ \hline
\texttt{outcome} & \mref{Verdict} & Outcome of this finding. \\ \hline
\texttt{basis} & \pref{String} & The rule the check enforces, e.g. \texttt{W1}, \texttt{V9}, \texttt{schema}. \\ \hline
\texttt{location} & \pref{String} & Where the finding applies, e.g. \texttt{step evaluate}. \\ \hline
\texttt{message} & \pref{String} & Human-readable explanation. \\ \hline
\end{longtable}
}

\msubsection{struct}{InterfaceEntry}
An interface of a service instance.

{\small
\begin{longtable}{|p{3.9cm}|p{4.4cm}|p{7.5cm}|}

\hline
\rowcolor{gray!33} Field & Type & Description \\ \hline
\endhead
\texttt{templateName} & \pref{String} & Interface template. \\ \hline
\texttt{policy} & \pref{String} & Interface policy, e.g. \texttt{CERT\_AUTH}. \\ \hline
\end{longtable}
}

\msubsection{enum}{LeaseMode}
Role of a step for a shared resource.

{\small
\begin{longtable}{|p{3.9cm}|p{11.9cm}|}

\hline
\rowcolor{gray!33} Value & Meaning \\ \hline
\endhead
\texttt{reserve} & Reserves the resource; its result carries the lease token. \\ \hline
\texttt{use} & Uses the resource; the executor renews the lease immediately before the invocation. \\ \hline
\texttt{release} & Releases the resource. \\ \hline
\end{longtable}
}

\msubsection{struct}{LeasePolicy}
The role of a step for a shared resource.

{\small
\begin{longtable}{|p{3.9cm}|p{4.4cm}|p{7.5cm}|}

\hline
\rowcolor{gray!33} Field & Type & Description \\ \hline
\endhead
\texttt{resource} & \pref{Name} & A declared resource. \\ \hline
\texttt{mode} & \mref{LeaseMode} & \texttt{reserve}, \texttt{use} or \texttt{release}. \\ \hline
\end{longtable}
}

\msubsection{enum}{OnError}
Reaction to an \texttt{ERROR} report.

{\small
\begin{longtable}{|p{3.9cm}|p{11.9cm}|}

\hline
\rowcolor{gray!33} Value & Meaning \\ \hline
\endhead
\texttt{rebind} & Exclude the executor for this session and bind the step again. \\ \hline
\texttt{retry} & Keep the executor: after \texttt{backoff\allowbreak{}Seconds} dispatch the step to the same executor again, or to the least-loaded candidate when that executor is no longer available (extension x-retry; for transient provider-side failures, the provider must accept a repeated call). \\ \hline
\texttt{abort} & Fail the session. \\ \hline
\end{longtable}
}

\msubsection{struct}{Parameter}
A declared session parameter.

{\small
\begin{longtable}{|p{3.9cm}|p{4.4cm}|p{7.5cm}|}

\hline
\rowcolor{gray!33} Field & Type & Description \\ \hline
\endhead
\texttt{name} & \pref{Name} & Parameter name, referenced as \texttt{\$\{param.name\}}. \\ \hline
\texttt{type} & \pref{String} & Informative type: \texttt{string}, \texttt{number}, \texttt{boolean} or \texttt{object}. Optional. \\ \hline
\texttt{required} & \pref{Boolean} & Whether \texttt{workflow\allowbreak{}Instance.\allowbreak{}start} refuses a session without it. Optional. \\ \hline
\end{longtable}
}

\msubsection{struct}{Policy}
Execution policy of a step. Every field other than \texttt{lease} and \texttt{local\allowbreak{}Cloud\allowbreak{}Only} is an extension outside the proven model.

{\small
\begin{longtable}{|p{3.9cm}|p{4.4cm}|p{7.5cm}|}

\hline
\rowcolor{gray!33} Field & Type & Description \\ \hline
\endhead
\texttt{quantity} & \pref{Integer} & Number of executions; only 1 is accepted (check M01). Optional. \\ \hline
\texttt{timeout\allowbreak{}Seconds} & \pref{Duration} & Extension x-timeout: the step fails when it runs longer. Optional. \\ \hline
\texttt{nominal\allowbreak{}Duration\allowbreak{}Seconds} & \pref{Duration} & Expected duration, used to check lease durations (check L04). Optional. \\ \hline
\texttt{retry} & \mref{Retry} & Extensions x-rebind and x-retry: behaviour on \texttt{ERROR}. Optional. \\ \hline
\texttt{lease} & \mref{LeasePolicy} & The step's role for a shared resource. Optional. \\ \hline
\texttt{local\allowbreak{}Cloud\allowbreak{}Only} & \pref{Boolean} & Whether orchestration may return providers of other clouds (default true). Optional. \\ \hline
\end{longtable}
}

\msubsection{struct}{Recipe}
A workflow recipe $W$~\cite{varga2026}: named, versioned and immutable once registered. Actions run in sequence (\texttt{firstAction}, then each action's \texttt{next}); the steps of one action form a directed acyclic graph.

{\small
\begin{longtable}{|p{3.9cm}|p{4.4cm}|p{7.5cm}|}

\hline
\rowcolor{gray!33} Field & Type & Description \\ \hline
\endhead
\texttt{name} & \pref{Name} & Recipe name. \\ \hline
\texttt{version} & \pref{Version} & Recipe version; a registered name and version never changes content. \\ \hline
\texttt{description} & \pref{String} & Free text. Optional. \\ \hline
\texttt{firstAction} & \pref{Name} & The action $\alpha_0$ that is active when a session starts. \\ \hline
\texttt{parameters} & \pref{List}$<$\mref{Parameter}$>$ & Declared session parameters $\rho$. Optional. \\ \hline
\texttt{resources} & \pref{List}$<$\mref{Resource}$>$ & Shared resources (tools) used under leases. Optional. \\ \hline
\texttt{dispatch} & \mref{DispatchMode} & Default dispatch mode of the recipe; overridden per action or per session. Optional. \\ \hline
\texttt{actions} & \pref{List}$<$\mref{Action}$>$ & The actions $A$ with their steps. \\ \hline
\end{longtable}
}

\msubsection{struct}{RecipeFilter}
Optional filter of the recipe query.

{\small
\begin{longtable}{|p{3.9cm}|p{4.4cm}|p{7.5cm}|}

\hline
\rowcolor{gray!33} Field & Type & Description \\ \hline
\endhead
\texttt{name} & \pref{Name} & Only versions of this recipe. Optional. \\ \hline
\texttt{summary} & \pref{Boolean} & Leave out the recipe documents (default false). Optional. \\ \hline
\end{longtable}
}

\msubsection{struct}{RecipeKey}
Identifies one registered recipe version.

{\small
\begin{longtable}{|p{3.9cm}|p{4.4cm}|p{7.5cm}|}

\hline
\rowcolor{gray!33} Field & Type & Description \\ \hline
\endhead
\texttt{name} & \pref{Name} & Recipe name. \\ \hline
\texttt{version} & \pref{Version} & Recipe version; for \texttt{compile} the latest active version when absent. \\ \hline
\end{longtable}
}

\msubsection{struct}{RecipeList}
Registered recipes in registration order.

{\small
\begin{longtable}{|p{3.9cm}|p{4.4cm}|p{7.5cm}|}

\hline
\rowcolor{gray!33} Field & Type & Description \\ \hline
\endhead
\texttt{entries} & \pref{List}$<$\mref{RecipeRecord}$>$ & The recipes. \\ \hline
\texttt{count} & \pref{Integer} & Number of entries. \\ \hline
\end{longtable}
}

\msubsection{struct}{RecipeRecord}
A registered recipe as stored by the WorkflowManager.

{\small
\begin{longtable}{|p{3.9cm}|p{4.4cm}|p{7.5cm}|}

\hline
\rowcolor{gray!33} Field & Type & Description \\ \hline
\endhead
\texttt{id} & \pref{Integer} & Database identifier. \\ \hline
\texttt{name} & \pref{Name} & Recipe name. \\ \hline
\texttt{version} & \pref{Version} & Recipe version. \\ \hline
\texttt{status} & \mref{RecipeStatus} & \texttt{ACTIVE} or \texttt{RETIRED}. \\ \hline
\texttt{registeredBy} & \pref{SystemName} & The system that registered the recipe. \\ \hline
\texttt{recipe} & \mref{Recipe} & The canonical recipe document; null in a \texttt{workflow\allowbreak{}Recipe.\allowbreak{}query} with \texttt{summary=true}. \\ \hline
\texttt{createdAt} & \pref{DateTime} & Registration time. \\ \hline
\end{longtable}
}

\msubsection{enum}{RecipeStatus}
Life-cycle state of a registered recipe.

{\small
\begin{longtable}{|p{3.9cm}|p{11.9cm}|}

\hline
\rowcolor{gray!33} Value & Meaning \\ \hline
\endhead
\texttt{ACTIVE} & Sessions may be started. \\ \hline
\texttt{RETIRED} & Kept for traceability; new sessions are refused with 409. \\ \hline
\end{longtable}
}

\msubsection{struct}{Resource}
A shared resource that steps reserve, use and release under a lease (Proposition~1 of~\cite{varga2026}, lease variant).

{\small
\begin{longtable}{|p{3.9cm}|p{4.4cm}|p{7.5cm}|}

\hline
\rowcolor{gray!33} Field & Type & Description \\ \hline
\endhead
\texttt{name} & \pref{Name} & Resource name. \\ \hline
\texttt{lease\allowbreak{}Duration\allowbreak{}Seconds} & \pref{Duration} & Lease duration granted by \texttt{workflowLease}; must exceed the longest tool-using step. Optional. \\ \hline
\texttt{provider} & \pref{SystemName} & Lease provider; the WorkflowManager when absent. Optional. \\ \hline
\end{longtable}
}

\msubsection{struct}{Retry}
Extensions x-rebind and x-retry.

{\small
\begin{longtable}{|p{3.9cm}|p{4.4cm}|p{7.5cm}|}

\hline
\rowcolor{gray!33} Field & Type & Description \\ \hline
\endhead
\texttt{maxAttempts} & \pref{Integer} & Maximum attempts of the step in one session, counting every work item of the step; at least 1, default 1 (no further attempt). Optional. \\ \hline
\texttt{onError} & \mref{OnError} & What an \texttt{ERROR} report causes while attempts remain. State it explicitly: the recipe schema names \texttt{rebind} as default, but the Java WorkflowManager fails the session when it is absent. Optional. \\ \hline
\texttt{backoff\allowbreak{}Seconds} & \pref{Duration} & Pause before a same-executor retry (\texttt{onError} = \texttt{retry}); 0--600, default 1. Optional. \\ \hline
\end{longtable}
}

\msubsection{struct}{ServiceEntry}
A service instance offered in the cloud.

{\small
\begin{longtable}{|p{3.9cm}|p{4.4cm}|p{7.5cm}|}

\hline
\rowcolor{gray!33} Field & Type & Description \\ \hline
\endhead
\texttt{provider} & \pref{SystemName} & Provider. \\ \hline
\texttt{service\allowbreak{}Definition} & \pref{ServiceDefinition} & Service. \\ \hline
\texttt{version} & \pref{String} & Service version. \\ \hline
\texttt{metadata} & \pref{JsonObject} & Service metadata. \\ \hline
\texttt{interfaces} & \pref{List}$<$\mref{InterfaceEntry}$>$ & Interfaces. \\ \hline
\texttt{operations} & \pref{List}$<$\pref{String}$>$ & Operation names. \\ \hline
\end{longtable}
}

\msubsection{struct}{ServiceList}
The microservices of a cloud.

{\small
\begin{longtable}{|p{3.9cm}|p{4.4cm}|p{7.5cm}|}

\hline
\rowcolor{gray!33} Field & Type & Description \\ \hline
\endhead
\texttt{cloud} & \pref{String} & Cloud identifier. \\ \hline
\texttt{name} & \pref{String} & Cloud name. \\ \hline
\texttt{live} & \pref{Boolean} & Live cloud. \\ \hline
\texttt{systems} & \pref{List}$<$\mref{CloudSystem}$>$ & Systems. \\ \hline
\texttt{entries} & \pref{List}$<$\mref{ServiceEntry}$>$ & Service instances. \\ \hline
\texttt{note} & \pref{String} & Why the list may be incomplete. Optional. \\ \hline
\end{longtable}
}

\msubsection{struct}{ServiceRequirement}
Requirements on the provider of the step's service, passed to orchestration.

{\small
\begin{longtable}{|p{3.9cm}|p{4.4cm}|p{7.5cm}|}

\hline
\rowcolor{gray!33} Field & Type & Description \\ \hline
\endhead
\texttt{service\allowbreak{}Definition} & \pref{ServiceDefinition} & The service to consume. \\ \hline
\texttt{interfaces} & \pref{List}$<$\pref{String}$>$ & Accepted interface template names, e.g. \texttt{generic\_\allowbreak{}https}. Optional. \\ \hline
\texttt{minVersion} & \pref{String} & Lowest accepted service version. Optional. \\ \hline
\texttt{maxVersion} & \pref{String} & Highest accepted service version. Optional. \\ \hline
\texttt{metadata} & \pref{Map}$<$\pref{String}, \pref{JsonValue}$>$ & Metadata requirements on the service instance. Optional. \\ \hline
\texttt{qos} & \pref{Map}$<$\pref{String}, \pref{JsonValue}$>$ & QoS requirements forwarded to orchestration. Optional. \\ \hline
\texttt{requires\allowbreak{}Results} & \pref{List}$<$\pref{Name}$>$ & Steps whose results the provider needs; each must be in $dep(s)$ (check D04). Optional. \\ \hline
\end{longtable}
}

\msubsection{struct}{SourceState}
Availability of a source.

{\small
\begin{longtable}{|p{3.9cm}|p{4.4cm}|p{7.5cm}|}

\hline
\rowcolor{gray!33} Field & Type & Description \\ \hline
\endhead
\texttt{system} & \pref{SystemName} & Source system. \\ \hline
\texttt{operation} & \pref{String} & Operation used. \\ \hline
\texttt{available} & \pref{Boolean} & Whether it answered. \\ \hline
\texttt{status} & \pref{Integer} & HTTP status when unavailable. Optional. \\ \hline
\texttt{error} & \pref{String} & Error when unavailable. Optional. \\ \hline
\end{longtable}
}

\msubsection{struct}{Step}
A step $s$: one service exchange.

{\small
\begin{longtable}{|p{3.9cm}|p{4.4cm}|p{7.5cm}|}

\hline
\rowcolor{gray!33} Field & Type & Description \\ \hline
\endhead
\texttt{name} & \pref{Name} & Step name, unique in the recipe. \\ \hline
\texttt{service} & \mref{ServiceRequirement} & $\theta(s)$: the service the step consumes. \\ \hline
\texttt{next} & \pref{List}$<$\pref{Name}$>$ & Successors in the same action (the transitions $s \rightarrow s'$). Optional. \\ \hline
\texttt{dep} & \pref{List}$<$\pref{Name}$>$ & $dep(s)$: steps whose results the step reads. Each must run before $s$. Optional. \\ \hline
\texttt{executor} & \mref{ExecutorRequirement} & $\eta(s)$: which executor may run the step; absent means any executor declaring the service. Optional. \\ \hline
\texttt{params} & \pref{Map}$<$\pref{String}, \pref{JsonValue}$>$ & Static parameters passed unchanged to the provider. Optional. \\ \hline
\texttt{input} & \pref{Template} & $\kappa_s$: the request body sent to the provider, composed from $\rho$ and $\xi|dep(s)$. Optional. \\ \hline
\texttt{policy} & \mref{Policy} & Execution policy (quantity, timeout, retry, lease, local cloud). Optional. \\ \hline
\texttt{guards} & \pref{Map}$<$\pref{Name}, \pref{Expression}$>$ & Conditions on the transitions to successors (extension x-guard). A successor without a guard is always taken. Optional. \\ \hline
\end{longtable}
}

\msubsection{struct}{StoreResult}
Result of storing an edited workflow.

{\small
\begin{longtable}{|p{3.9cm}|p{4.4cm}|p{7.5cm}|}

\hline
\rowcolor{gray!33} Field & Type & Description \\ \hline
\endhead
\texttt{stored} & \pref{Boolean} & True when the WorkflowManager registered it. \\ \hline
\texttt{report} & \mref{ValidationReport} & The validation report. \\ \hline
\texttt{created} & \pref{Boolean} & True for a new version; false for identical content. Optional. \\ \hline
\texttt{recipe} & \mref{RecipeRecord} & The stored recipe. Optional. \\ \hline
\end{longtable}
}

\msubsection{struct}{ValidationReport}
The result of validating a recipe. The verdict is the worst outcome of the findings (REJECT $>$ CLARIFY $>$ REVISE $>$ ACCEPT).

{\small
\begin{longtable}{|p{3.9cm}|p{4.4cm}|p{7.5cm}|}

\hline
\rowcolor{gray!33} Field & Type & Description \\ \hline
\endhead
\texttt{recipe} & \pref{String} & \texttt{name@version} of the validated recipe. \\ \hline
\texttt{verdict} & \mref{Verdict} & Overall verdict; only \texttt{ACCEPT} can be registered. \\ \hline
\texttt{findings} & \pref{List}$<$\mref{Finding}$>$ & Individual findings; empty when the recipe has none. \\ \hline
\end{longtable}
}

\msubsection{enum}{Verdict}
Validator outcome.

{\small
\begin{longtable}{|p{3.9cm}|p{11.9cm}|}

\hline
\rowcolor{gray!33} Value & Meaning \\ \hline
\endhead
\texttt{ACCEPT} & The recipe can be registered (also used for informational findings). \\ \hline
\texttt{REVISE} & The recipe must be changed. \\ \hline
\texttt{CLARIFY} & The recipe refers to something that cannot be confirmed (undeclared parameter, unknown service or executor). \\ \hline
\texttt{REJECT} & The recipe cannot be used at all. \\ \hline
\end{longtable}
}


\subsection{Primitives}
\label{sec:model:primitives}
Types and structures mentioned throughout this document that are assumed to be available to implementations of this
service. The concrete interpretation of each of these types must be provided by any IDD document claiming to implement
this service.

{\small
\begin{longtable}{|p{3.2cm}|p{12.6cm}|}

\hline
\rowcolor{gray!33} Type & Description \\ \hline
\endhead
\pdef{Boolean} & One out of \texttt{true} or \texttt{false}. \\ \hline
\pdef{DateTime} & Pinpoints a moment in time (UTC). \\ \hline
\pdef{Duration} & A length of time in seconds. \\ \hline
\pdef{Expression} & A side-effect free condition over the session parameters and stored results (Section on guard expressions). \\ \hline
\pdef{Integer} & A whole number. \\ \hline
\pdef{JsonObject} & A set of named values. \\ \hline
\pdef{JsonValue} & Any value: object, array, string, number, Boolean or null. \\ \hline
\pdef{List} & An ordered collection of items of the type given in angle brackets. \\ \hline
\pdef{Map} & A collection of key/value pairs of the types given in angle brackets. \\ \hline
\pdef{Name} & An identifier of a recipe, action, step, parameter or resource: a letter followed by at most 62 letters, digits, \texttt{-} or \texttt{\_}. \\ \hline
\pdef{ServiceDefinition} & The name of an Arrowhead service definition. Names are compared case-insensitively and without \texttt{-}, \texttt{\_} and spaces, so \texttt{temperature-\allowbreak{}read} and \texttt{temperature\allowbreak{}Read} denote the same service. \\ \hline
\pdef{String} & A chain of characters. \\ \hline
\pdef{SystemName} & The name of an Arrowhead system as it appears in its certificate and in the ServiceRegistry. \\ \hline
\pdef{Template} & A value in which strings of the form \texttt{\$\{param.$p$\}} or \texttt{\$\{result.$s$.path\}} are replaced by a session parameter or by (a part of) a stored step result. \\ \hline
\pdef{Version} & A recipe version in the form MAJOR.MINOR.PATCH. \\ \hline
\end{longtable}
}


\newpage
\bibliographystyle{IEEEtran}
\bibliography{bibliography}

\newpage

\section{Revision History}
\subsection{Amendments}

\noindent\begin{tabularx}{\textwidth}{| p{1cm} | p{2.2cm} | p{1.6cm} | X | p{3.6cm} |} \hline
\rowcolor{gray!33} No. & Date & Version & Subject of Amendments & Author \\ \hline
1 & 2026-11-23 & \arrowversion & Initial edition: service description of workflowConfiguration, written from the Java reference implementation (ah5-Workflow-java). & Documentation maintenance \\ \hline
2 & 2026-11-24 & \arrowversion & Retry policy as implemented. & Documentation maintenance \\ \hline
\end{tabularx}

\subsection{Quality Assurance}

\noindent\begin{tabularx}{\textwidth}{| p{1cm} | p{2.2cm} | p{1.6cm} | X |} \hline
\rowcolor{gray!33} No. & Date & Version & Approved by \\ \hline
1 & & \arrowversion & Pending review \\ \hline
\end{tabularx}

\end{document}
